% TOP_PROBLEM: {"problemId": "problem.global-gan-gross-prasad-conjecture", "canonicalTitle": "Global Gan–Gross–Prasad Conjecture (Fixed-Representation Form)", "releaseRank": 138, "primaryDomain": "number_theory_arithmetic_geometry", "primaryDomainLabel": "Number theory & arithmetic geometry", "displayStatus": "open_with_solved_subcases", "releaseStatus": "open_with_solved_subcases"}
% !TeX program = lualatex
% Definition and sources only; this file is not an LLM solution attempt.
\documentclass[11pt]{article}
\usepackage[margin=25mm]{geometry}
\usepackage{fontspec}
\setmainfont{Latin Modern Roman}
\usepackage{amsmath,amssymb,mathtools}
\usepackage{unicode-math}
\setmathfont{Latin Modern Math}
\usepackage{xurl,hyperref}
\hypersetup{colorlinks=true,urlcolor=blue}
\setcounter{secnumdepth}{0}
\setlength{\parindent}{0pt}
\setlength{\parskip}{5pt}
\setlength{\emergencystretch}{3em}
\begin{document}
\section*{\texorpdfstring{Global Gan–Gross–Prasad Conjecture (Fixed-Representation Form)}{Global Gan–Gross–Prasad Conjecture (Fixed-Representation Form)}}\label{138}
\subsection{Definitions and mathematical statement}
For the relevant classical-group pair, let $H$ be the subgroup in Gan--Gross--Prasad, $\nu=\bigotimes'_v\nu_v$ its Bessel or Fourier--Jacobi data, and $R$ the prescribed representation defining the finite $L$-function. For each fixed tempered cuspidal automorphic representation $\pi$ occurring with multiplicity one, the assertion is
\[
 \mathcal F(\nu)|_\pi\ne0
 \quad\Longleftrightarrow\quad
 \left(\forall v,\ \operatorname{Hom}_{H(F_v)}(\pi_v,\nu_v)\ne0\right)
 \ \text{and}\ L(\pi,R,\tfrac12)\ne0.
\]
The left side says that the source-defined global period is not identically zero on $\pi$. The local Hom-spaces test distinction at every place. The form signs, covers, auxiliary characters, and inverse twist in $R$ are retained from Sections 23--24. Neither a packet-existence statement nor a refined numerical period formula is substituted.
\par\smallskip\textit{Formulation and concept references:} \href{https://arxiv.org/pdf/0909.2999}{[S1]}.

\subsection{Short English statement}
For each allowed automorphic representation, is its global period nonzero exactly when every local period is allowed and the central $L$-value is nonzero?
\subsection{Sources}
\begin{itemize}
\item[S1]Gan, Gross and Prasad, Symplectic local root numbers, central critical L-values, and restriction problems in the representation theory of classical groups, §§23–24, Conjecture 24.1.
\url{https://arxiv.org/pdf/0909.2999}.
\textit{Formulation source.}
\end{itemize}
\end{document}
